\section{Spectral exponential weights and truncation excesses}

All matrix norms in this section are Euclidean operator norms. These estimates
support the regional truncation argument
in~\cite[lines 1102--1168]{Arad2014Connecting} and
\cite[lines 1010--1024]{Arad2016RigorousRG}.

\begin{theorem}[Exponential weights and their spectral projections]
    \label{thm:exponential_weight_projection}
    \lean{Matrix.IsHermitian.regionalWeight_relations,
        Matrix.IsHermitian.norm_exp_mul_sub_const_mul_lowProjection_le_one}
    \leanok
    \uses{def:closed_threshold_projection}
    Let $A$ be Hermitian and $s,c\in\R$. Put
    $W=e^{s(A-c\Id)}$, $W_-=e^{-s(A-c\Id)}$ and $Q=\mathbf1_{\{x\ge c\}}(A)$.
    Then $WQ=QW$ and $W_-W=\Id$. Any matrix $X$ commuting with $A$ commutes
    with $W$ and $Q$. If $s\ge0$, then $\|W(1-Q)\|\le1$.
\end{theorem}
\begin{proof}
    \leanok
    All the stated weights and projections are functions of $A$.
    On its spectrum, their scalar products give the commutations and inverse.
    For $s\ge0$, the low part has scalar value $e^{s(r-c)}$ when $r<c$
    and zero otherwise, whose absolute value is at most one.
\end{proof}

\begin{theorem}[Norms of projected and forcing terms]
    \label{thm:weighted_projection_norms}
    \lean{Matrix.norm_projection_conjugated_mul_projection_le,
        Matrix.norm_weighted_projection_forcing_le}
    \leanok
    Let $Q$ be an orthogonal projection and $W,Y,W_-$ arbitrary square
    complex matrices. Then $\|QWYW_-Q\|\le\|WYW_-\|$.
    If $WQ=QW$ and $W_-W=\Id$, then also
    \begin{align}
        \|WQY(1-Q)\|\le\|WYW_-\|\,\|W(1-Q)\|.
        \label{eq:weighted_forcing_norm}
    \end{align}
\end{theorem}
\begin{proof}
    \leanok
    \uses{thm:weighted_forcing_identity}
    Orthogonal projections are contractions. This gives the first estimate.
    For the second, use~(\ref{eq:weighted_forcing_factorization}) with $R=\Id$,
    and apply submultiplicativity and $\|Q\|\le1$.
\end{proof}

\begin{theorem}[A linear factor under an exponential weight]
    \label{thm:linear_exponential_weight}
    \lean{Real.sub_mul_exp_weight_le, Real.posPart_sub_mul_exp_weight_le}
    \leanok
    For $s>0$ and arbitrary $x,t,u\in\R$,
    \begin{align}
        (x-t)e^{-s(x-u)}\le\frac{e^{-s(t-u)-1}}s.
        \label{eq:linear_exponential_weight}
    \end{align}
    The same estimate holds with $x-t$ replaced by $\max(x-t,0)$.
\end{theorem}
\begin{proof}
    \leanok
    Apply $ye^{-y}\le e^{-1}$ to $y=s(x-t)$, multiply by
    $e^{-s(t-u)}$ and divide by $s$. The right side is positive, so taking
    the positive part of the left factor preserves the bound.
\end{proof}

\begin{theorem}[Truncation excess on an actual weighted restriction]
    \label{thm:excess_actual_weighted_restriction}
    \lean{Matrix.IsHermitian.norm_excess_domRestrict_le_actual_weighted_map}
    \leanok
    \uses{def:energy_truncation, def:closed_threshold_projection}
    Let $A$ be Hermitian, $s>0$ and $u\le t$. Set
    $W=e^{s(A-u\Id)}$ and $Q=\mathbf1_{\{x\ge u\}}(A)$.
    Let $G,U$ be arbitrary subspaces of the Euclidean space. If a linear
    map $Z:G\to U$ satisfies $\iota_UZ=(WQ)|_G$, then
    \begin{align}
        \|(A-\min(A,t))|_G\|
        \le\frac{e^{-s(t-u)-1}}s\,\|Z\|.
        \label{eq:excess_weighted_restriction}
    \end{align}
    In particular, no norm estimate for $Z$ is assumed in this assertion.
\end{theorem}
\begin{proof}
    \leanok
    \uses{thm:linear_exponential_weight, thm:weighted_right_support_cancellation}
    Put $D=A-\min(A,t)$ and $W_-=e^{-s(A-u\Id)}$.
    Since $u\le t$, functional calculus gives $D(1-Q)=0$.
    The scalar estimate bounds $\|DW_-\|$ by $e^{-s(t-u)-1}/s$.
    The cancellation identity gives $(DW_-)(WQ)=D$.
    Restrict to $G$, substitute the actual embedding of $Z$, and take norms.
\end{proof}
